\chapter{Procesor jednocyklowy RV32I}
\label{ch:single}

\metryczka{złożyć działający procesor, który wykonuje prawdziwe programy w kodzie
maszynowym RISC-V, po jednej instrukcji na takt zegara.}{3--4 wieczory}{05-single-cycle/}

\section{Ścieżka danych}

\begin{figure}[h]
\centering
\resizebox{\linewidth}{!}{% Ścieżka danych procesora jednocyklowego (rozdział 5)
\begin{tikzpicture}[x=1cm, y=1cm, every node/.append style={font=\sffamily\scriptsize},
  szyna/.style={thick}, ster/.style={densely dashed, pulapka, thin, ->}]
  % --- bloki
  \node[draw, thick, fill=zadaniejasny, minimum width=6mm, minimum height=16mm] (pc) at (0,0) {PC};
  \node[blok, minimum width=15mm, minimum height=16mm] (imem) at (1.9,0) {pamięć\\instrukcji};
  \node[blok, minimum width=15mm, minimum height=9mm] (ctrl) at (5.4,2.75) {control};
  \node[blok, minimum width=16mm, minimum height=20mm] (rf) at (5.4,0) {regfile};
  \node[blok, minimum width=16mm, minimum height=8mm] (imm) at (5.4,-2.1) {imm\_gen};
  \node[blok, minimum width=7mm, minimum height=10mm] (ma) at (7.9,0.55) {A};
  \node[blok, minimum width=7mm, minimum height=10mm] (mb) at (7.9,-0.55) {B};
  \node[blok, minimum width=12mm, minimum height=24mm, fill=pytaniejasny] (alu) at (9.5,0) {ALU};
  \node[blok, minimum width=17mm, minimum height=8mm, fill=pytaniejasny] (br) at (9.5,-2.7) {branch\_cmp};
  \node[blok, minimum width=17mm, minimum height=16mm] (mem) at (11.9,0) {LSU +\\pamięć\\danych};
  \node[blok, minimum width=7mm, minimum height=14mm] (wb) at (13.9,0) {WB};
  \node[blok, minimum width=15mm, minimum height=9mm] (npc) at (1.9,-2.6) {następny\\PC};
  % --- PC -> imem -> instr
  \draw[szyna, ->] (pc) -- (imem);
  \draw[szyna] (imem.east) -- (3.4,0);
  \draw[szyna] (3.4,2.75) -- (3.4,-2.1);
  \node[sygnal, rotate=90, fill=white, inner sep=1pt] at (3.4,1.4) {instr};
  \draw[szyna, ->] (3.4,2.75) -- (ctrl.west);
  \draw[szyna, ->] (3.4,0.45) -- node[above, sygnal]{rs1, rs2, rd} ($(rf.west)+(0,0.45)$);
  \draw[szyna, ->] (3.4,-2.1) -- (imm.west);
  % --- regfile -> A/B -> ALU
  \draw[szyna, ->] ($(rf.east)+(0,0.55)$) -- node[above, sygnal, pos=0.75]{rs1} (ma.west);
  \draw[szyna, ->] ($(rf.east)+(0,-0.55)$) -- node[below, sygnal, pos=0.75]{rs2} (mb.west);
  \draw[szyna, ->] (imm.east) -- (7.9,-2.1) node[right, sygnal, pos=0.4, yshift=5pt]{} -- (mb.south);
  \node[sygnal] at (8.15,-1.6) {imm};
  \draw[szyna, <-] (ma.north) -- ++(0,0.45) node[left, sygnal]{pc, 0};
  \draw[szyna, ->] (ma.east) -- ($(alu.west)+(0,0.55)$);
  \draw[szyna, ->] (mb.east) -- ($(alu.west)+(0,-0.55)$);
  % --- komparator (odczepy rs1, rs2)
  \fill (6.55,0.55) circle (1.3pt);
  \fill (6.85,-0.55) circle (1.3pt);
  \draw[szyna, ->] (6.55,0.55) -- (6.55,-2.55) -- ($(br.west)+(0,0.15)$);
  \draw[szyna, ->] (6.85,-0.55) -- (6.85,-2.85) -- ($(br.west)+(0,-0.15)$);
  % --- ALU -> pamięć -> WB
  \draw[szyna, ->] (alu.east) -- node[above, sygnal]{alu\_y} (mem.west);
  \draw[szyna, ->] (mem.east) -- node[above, sygnal, pos=0.45]{load} (wb.west);
  \fill (10.55,0) circle (1.3pt);
  \draw[szyna, ->] (10.55,0) -- (10.55,-1.3) -| ($(wb.south)+(-0.15,0)$);
  \draw[szyna, <-] ($(wb.south)+(0.15,0)$) -- ++(0,-0.8) node[below, sygnal]{pc+4};
  % --- zapis z WB do regfile
  \draw[szyna, ->] (wb.east) -- ++(0.35,0) |- (5.9,2.2) -- ($(rf.north)+(0.5,0)$);
  \node[sygnal, fill=white, inner sep=1pt] at (11.2,2.2) {wynik do rejestru rd};
  % --- następny PC: imm, taken, alu_y
  \draw[szyna, ->] (imm.south) -- (5.4,-3.55) -| ($(npc.south)+(-0.35,0)$);
  \node[sygnal, fill=white, inner sep=1pt] at (4.0,-3.55) {imm};
  \draw[szyna, ->] (br.south) -- (9.5,-3.85) -| (npc.south);
  \node[sygnal, fill=white, inner sep=1pt] at (7.6,-3.85) {taken};
  \fill (10.55,-1.3) circle (1.3pt);
  \draw[szyna, ->] (10.55,-1.3) -- (10.55,-4.15) -| ($(npc.south)+(0.35,0)$);
  \node[sygnal, fill=white, inner sep=1pt] at (8.5,-4.15) {alu\_y (jalr)};
  \draw[szyna, ->] (npc.west) -- (-0.75,-2.6) |- (pc.west);
  \node[sygnal, rotate=90] at (-1.0,-1.3) {pc+4 / pc+imm / alu\_y};
  % --- sterowanie (przerywane)
  \draw[ster] (ctrl.east) -| ($(ma.north)+(0.3,0)$);
  \draw[ster] ($(ctrl.east)+(0,0.15)$) -| (alu.north);
  \draw[ster] ($(ctrl.east)+(0,0.3)$) -| (wb.north);
  \node[sygnal, text=pulapka, anchor=west] at (6.3,3.25) {sygnały sterujące (wb\_sel, alu\_op, \dots)};
\end{tikzpicture}
}
\caption{Ścieżka danych procesora jednocyklowego. Linie przerywane to sygnały
sterujące z dekodera.}
\label{fig:datapath}
\end{figure}

W jednym takcie, od zbocza do zbocza (rys.~\ref{fig:datapath}):
\begin{enumerate}
  \item \code{pc} adresuje pamięć instrukcji, a \code{instr} pojawia się kombinacyjnie.
  \item Dekoder (\code{control}) i \code{imm\_gen} analizują \code{instr}, a bank
    rejestrów czyta \code{rs1} i \code{rs2}.
  \item ALU liczy wynik, adres dla load/store albo cel skoku dla \code{jalr}.
  \item Load: LSU wybiera bajty ze słowa z pamięci danych. Store: LSU ustawia
    \code{wstrb}/\code{wdata}, a zapis nastąpi na zboczu.
  \item Multiplekser WB wybiera, co trafi do \code{rd}: wynik ALU, daną z
    pamięci albo \code{pc+4}. Zapis nastąpi na zboczu.
  \item Logika następnego PC wybiera \code{pc+4}, \code{pc+imm} albo
    \code{alu\_y \& \textasciitilde1}. Nowa wartość też trafia do PC na zboczu.
\end{enumerate}

\textbf{Wszystkie efekty instrukcji (zapis rejestru, zapis pamięci, nowe PC)
dzieją się na tym samym zboczu zegara.} Dlatego ten procesor jest taki prosty.
Cena: okres zegara musi pomieścić całą drogę od PC przez pamięć instrukcji,
rejestry, ALU i pamięć danych aż do zapisu.

\section{Przejście instrukcji przez ścieżkę danych}

\begin{center}
\small
\begin{tabular}{@{}llllll@{}}
\toprule
\textbf{Instrukcja} & \textbf{A} & \textbf{B} & \textbf{ALU} & \textbf{Do rd} & \textbf{Następne PC} \\
\midrule
\code{add rd, rs1, rs2} & rs1 & rs2 & \{i[30], f3\} & alu\_y & pc+4 \\
\code{addi rd, rs1, imm} & rs1 & imm & \{0, f3\}$^*$ & alu\_y & pc+4 \\
\code{lw rd, imm(rs1)} & rs1 & imm & ADD & load\_data & pc+4 \\
\code{sw rs2, imm(rs1)} & rs1 & imm & ADD & --- (zapis rs2) & pc+4 \\
\code{beq rs1, rs2, off} & --- & --- & --- & --- & taken ? pc+imm : pc+4 \\
\code{jal rd, off} & --- & --- & --- & pc+4 & pc+imm \\
\code{jalr rd, imm(rs1)} & rs1 & imm & ADD & pc+4 & alu\_y \& \textasciitilde1 \\
\code{lui rd, imm} & 0 & imm & ADD & alu\_y & pc+4 \\
\code{auipc rd, imm} & pc & imm & ADD & alu\_y & pc+4 \\
\bottomrule
\end{tabular}

\smallskip
{\footnotesize $^*$ dla \code{srai}: \{i[30], f3\}}
\end{center}

\section{Tabela sterowania}

Wypełnij tę tabelę \textbf{na kartce}, zanim napiszesz \code{control.sv}.
Kreska oznacza „obojętne”. Test \code{tb\_control} sprawdza dokładnie tę tabelę na 97
zakodowanych instrukcjach i wypisuje różnice z nazwami pól.

\begin{center}
\small
\renewcommand{\arraystretch}{1.6}
\begin{tabular}{|l|p{9mm}|p{9mm}|p{9mm}|p{9mm}|p{9mm}|p{9mm}|p{9mm}|p{9mm}|p{9mm}|p{9mm}|}
\hline
& \rotatebox{90}{\code{reg\_write}\ } & \rotatebox{90}{\code{wb\_sel}\ } & \rotatebox{90}{\code{alu\_a\_sel}\ } & \rotatebox{90}{\code{alu\_b\_imm}\ } & \rotatebox{90}{\code{alu\_op}\ } & \rotatebox{90}{\code{mem\_read}\ } & \rotatebox{90}{\code{mem\_write}\ } & \rotatebox{90}{\code{branch}\ } & \rotatebox{90}{\code{jump}\ } & \rotatebox{90}{\code{jump\_reg}\ } \\
\hline
\textsf{OP (R)} & & & & & & & & & & \\
\hline
\textsf{OP-IMM} & & & & & & & & & & \\
\hline
\textsf{LOAD} & & & & & & & & & & \\
\hline
\textsf{STORE} & & & & & & & & & & \\
\hline
\textsf{BRANCH} & & & & & & & & & & \\
\hline
\textsf{JAL} & & & & & & & & & & \\
\hline
\textsf{JALR} & & & & & & & & & & \\
\hline
\textsf{LUI} & & & & & & & & & & \\
\hline
\textsf{AUIPC} & & & & & & & & & & \\
\hline
\textsf{FENCE/SYSTEM} & & & & & & & & & & \\
\hline
\end{tabular}
\end{center}

\begin{pulapka}[\texttt{addi} z ujemną stałą to nie \texttt{sub}]
Dla OP-IMM \code{instr[30]} jest częścią immediate'a. Gdyby dekoder brał go do
\code{alu\_op} jak dla R-type, \code{addi a0, a0, -1} (stała z ustawionym bitem
30) liczyłoby się jako odejmowanie. Bit 30 ma znaczenie tylko dla \code{srai}
(\code{funct3 = 101}).
\end{pulapka}

\section{LSU: bajty w słowie}

Pamięć jest zorganizowana w 32-bitowe słowa. \code{addr[31:2]} wybiera słowo, a
\code{addr[1:0]} bajt w słowie. Zapis częściowy realizuje maska
\code{wstrb} (\emph{write strobe}): bit $i$ = 1 oznacza zapis bajtu $i$.

\begin{figure}[h]
\centering
\begin{tikzpicture}[x=15mm]
  \foreach \i/\r in {0/{[7:0]},1/{[15:8]},2/{[23:16]},3/{[31:24]}} {
    \draw[thick, fill=akcentjasny] (3-\i,0) rectangle ++(1,0.7);
    \node[sygnal] at (3.5-\i,0.35) {bajt \i};
    \node[sygnal, text=szary] at (3.5-\i,-0.25) {\r};
  }
  \draw[thick, fill=pulapkajasny] (1,0) rectangle ++(1,0.7);
  \node[sygnal] at (1.5,0.35) {bajt 2};
  \node[anchor=west, font=\sffamily\small] at (4.4,0.35) {\code{sb x5, 2(x0)}: \code{wstrb = 4'b0100}, \code{wdata[23:16] = x5[7:0]}};
\end{tikzpicture}
\caption{Zapis bajtu pod adres o \code{addr[1:0] = 2}.}
\end{figure}

Przy odczycie pamięć zwraca całe słowo, a LSU przesuwa je o
\code{8*addr[1:0]} i rozszerza znakiem (\code{lb}, \code{lh}) albo zerami
(\code{lbu}, \code{lhu}).

\section{Interfejs rdzenia: kontrakt z testbenchem}

\begin{sv}
module core import rv32i_pkg::*; (
  input  logic        clk, rst,
  output logic [31:0] imem_addr,  input logic [31:0] imem_rdata,
  output logic [31:0] dmem_addr,  output logic dmem_re,
  input  logic [31:0] dmem_rdata,
  output logic [3:0]  dmem_wstrb, output logic [31:0] dmem_wdata,
  output logic        commit_valid,                    // ślad wykonania
  output logic [31:0] commit_pc, commit_instr, commit_rd_val,
  output logic [4:0]  commit_rd
);
\end{sv}

\begin{itemize}
  \item Pamięć (64 KiB) jest \emph{jedna}, widziana przez dwa porty. Odczyt jest
    kombinacyjny, a zapis następuje na zboczu, gdy \code{dmem\_wstrb != 0}.
  \item \code{dmem\_addr} to pełny adres bajtowy, a \code{dmem\_rdata} całe
    słowo. Wybór bajtów robi Twoje LSU.
  \item Reset synchroniczny, a po nim \code{pc = 0}.
  \item \textbf{Interfejs commit} służy tylko do weryfikacji. W każdym takcie,
    w którym kończy się instrukcja, \code{commit\_valid = 1} i opis tej
    instrukcji. \code{commit\_rd = 0}, gdy instrukcja nie zapisuje rejestru.
\end{itemize}

Testbench (\plik{common/tb/soc\_tb.sv}) przerywa symulację z czytelnym
komunikatem, gdy: PC jest X po resecie, \code{commit\_valid} jest X, następuje
zapis pod adres X albo nieobsługiwany adres, wykona się słowo
\code{0x00000000} (skok w pustą pamięć), przez 1000 taktów nie ma commitu albo
przekroczony zostanie limit cykli.

\section{Jak debugować procesor}

\begin{enumerate}
  \item Zacznij od \textbf{najprostszego testu}, który nie przechodzi:
    \code{make prog P=01\_smoke}.
  \item \textbf{Ślad.} \code{make prog P=...} zapisuje \plik{build/P.trace}.
    Porównaj go z emulatorem:
\begin{sh}
make -C ../04-isa-emulator emu-test P=05_branch
python3 ../common/tools/tracecmp.py ../04-isa-emulator/build/05_branch.trace \
                                    build/05_branch.trace
\end{sh}
    Wynik pokazuje pierwszą różniącą się instrukcję z disasemblacją i
    podpowiedzią: inne PC oznacza błąd skoku, inna wartość błąd wykonania.
  \item \textbf{Listing} \plik{build/hex/P.lst} łączy adresy z liniami źródła.
  \item \textbf{Przebiegi}: \code{make wave P=05\_branch}. Dodaj \code{dut.pc},
    \code{imem\_rdata}, sygnały sterujące i \code{commit\_*}.
  \item „FAIL: podtest N” oznacza, że w źródle testu przed nieudanym sprawdzeniem
    stoi \code{li gp, N}.
\end{enumerate}

Przykładowy raport \code{tracecmp} dla błędu w dekoderze:

\begin{txt}
RÓŻNICA w instrukcji nr 12:
  ostatnie zgodne:
    00000024: 00500593  addi a1, zero, 5                 x11=00000005
  oczekiwano (ref): 00000028: fff58593  addi a1, a1, -1   x11=00000004
  jest       (RTL): 00000028: fff58593  addi a1, a1, -1   x11=00000006
  -> ta sama instrukcja, inny efekt: błąd wykonania (ALU, LSU, forwarding...)
\end{txt}

\section{Programy testowe}

\begin{center}
\small
\begin{tabularx}{\linewidth}{@{}l X@{}}
\toprule
\code{01\_smoke} & tylko \code{lui}, \code{addi}, \code{sw}: najprostszy możliwy test \\
\code{02\_alu\_imm}, \code{03\_alu\_reg} & ALU z immediate i rejestr-rejestr, przypadki brzegowe \\
\code{04\_lui\_auipc} & stałe górne i adresowanie względem PC \\
\code{05\_branch} & wszystkie skoki warunkowe, także dalekie ($>$2 KiB) \\
\code{06\_jal\_jalr} & adresy powrotu, \code{rd == rs1}, stos, zagnieżdżone wywołania \\
\code{07\_load\_store} & wszystkie szerokości, rozszerzanie znakiem i zerami \\
\code{08\_fib}, \code{09\_sort}, \code{10\_recursion} & „prawdziwe” programy: pętle, tablice, rekurencja \\
\code{11\_hazards} & sekwencje groźne dla potoku (rozdział 8) \\
\code{12\_mmio} & wypisuje napis przez PUTCHAR i czyta licznik cykli \\
\bottomrule
\end{tabularx}
\end{center}

Gdy przejdzie \code{12\_mmio}, zobaczysz w terminalu tekst wypisany przez
Twój procesor.

\section{Zadania}

\begin{zadanie}{5.1 --- Dekoder (\code{rtl/control.sv})}
\code{make unit T=control}.
\end{zadanie}

\begin{zadanie}{5.2 --- LSU (\code{rtl/lsu.sv})}
\code{make unit T=lsu}. Rozrysuj sobie słowo z czterema bajtami i zobacz, gdzie
trafia \code{sb} pod \code{addr[1:0] = 3}.
\end{zadanie}

\begin{zadanie}{5.3 --- Rdzeń (\code{rtl/core.sv})}
Złóż ścieżkę danych z gotowych bloków. Moduły z rozdziałów 3 i 4 są dołączane z
tamtych katalogów, więc poprawka w \plik{alu.sv} od razu działa także tutaj.
\code{make progs} uruchamia 12 programów, \code{make prog P=08\_fib} jeden, a
\code{make test} wszystko.
\end{zadanie}

\begin{zadanie}{5.4 --- Pomiary}
\code{make synth T=core}: ile komórek ma Twój procesor w porównaniu z sumą ALU
i banku rejestrów? CPI wynosi z definicji 1,000. W rozdziale 8 porównasz to z
potokiem.
\end{zadanie}

\begin{uwaga}[Uwaga o długości ścieżki]
Pamięci są w testbenchu, poza modułem \code{core}, więc Yosys widzi osobne
kawałki: PC → \code{imem\_addr}, \code{imem\_rdata} → … → \code{dmem\_addr} oraz
\code{dmem\_rdata} → rejestry. Prawdziwy takt procesora jednocyklowego to
\emph{suma}: odczyt imem + dekodowanie + rejestry + ALU + odczyt dmem + zapis.
To jest motywacja potoku.
\end{uwaga}

\begin{gotowe}
\code{make test} → 2 × PASS (testy jednostkowe) + 12 × PASS (programy), \code{make lint} czysto.
\end{gotowe}

\begin{kontrolne}
  \item Wymień wszystkie elementy, przez które przechodzi \code{lw} w ciągu jednego taktu.
  \item Po co multiplekser A ma wejścia \code{pc} i \code{0}? Które instrukcje ich używają?
  \item Dlaczego cel \code{jalr} liczy ALU, a cel \code{jal} i \code{beq} osobny sumator?
  \item Jakie \code{wstrb} i \code{wdata} wystawi LSU dla \code{sh x7, 2(x0)}, gdy \code{x7 = 0x1234BEEF}?
  \item Dlaczego procesor jednocyklowy ma niską maksymalną częstotliwość zegara?
\end{kontrolne}
